\chapter{Kelly sizing}\label{ch:kelly}

\noindent\emph{This chapter is mathematics, not advice. Nothing here is a recommendation
to trade, size, or bet on anything; it describes a formula, its assumptions, and where
each fails.}

\vspace{1em}
\noindent Given an edge, how much of the bankroll should each trade risk? The Kelly
criterion answers with the fraction that maximizes the expected log of wealth---and it
is also the fraction that maximizes ruin when the edge is misestimated. Both facts are
the chapter.

\section{The formula}\label{sec:kelly}

For a bet paying $b$ per unit risked with win probability $p$, the growth rate per trial
at fraction $f$ is
\[
g(f)=p\log(1+fb)+(1-p)\log(1-f),
\]
maximized at $f^\star=(bp-(1-p))/b$. \texttt{kelly\_sims.py}:
\begin{lstlisting}[numbers=none]
coin p=0.55 b=1.0:  f* = 0.10   (theory 0.10)
growth at f* +0.0050/toss   at 2f* -0.0001   at 3f* -0.0162
f=0.10: final   61089.0   median path dd 97.8%
f=0.05: final      40.5   median path dd 92.1%
f=0.20: final      98.5   median path dd 100.0%
\end{lstlisting}

\begin{example}\label{ex:kelly}
\textbf{What f* actually optimizes.} On a 55/45 even-money coin, $f^\star=0.10$ grows at
$+0.5\%$ per toss---and the median path still draws down $97.8\%$.
\begin{solution}
Kelly maximizes asymptotic growth, not comfort. Doubling the bet ($f=0.20$) destroys
growth ($\approx0$) while tripling is negative---the curve is asymmetric: under-betting
costs growth, over-betting costs everything. The final wealth column is the point:
full Kelly $61{,}089\times$ vs double Kelly $98.5\times$ on the same edge.
\end{solution}
\end{example}

\begin{figure}[htbp]\centering
\begin{tikzpicture}
\begin{axis}[width=0.86\linewidth,height=0.36\linewidth,
    xlabel={fraction $f$},ylabel={growth/toss},
    xmin=0,xmax=0.35,
    tick style={draw=none},axis line style={draw=slate},font=\scriptsize]
  \addplot[accent,very thick] file{figures/data/KG.tex};
  \addplot[red,dashed,thick] coordinates {(0.1,0)(0.1,0.0055)};
  \addplot[slate,dotted] coordinates {(0,-0.02)(0.35,-0.02)};
  \node[font=\scriptsize,text=red] at (axis cs:0.115,0.0058) {$f^*{=}0.10$};
  \node[font=\scriptsize,text=slate] at (axis cs:0.26,-0.006) {$2f^*$: growth $\approx0$};
  \node[font=\scriptsize,text=slate] at (axis cs:0.30,-0.014) {$3f^*$: negative};
\end{axis}
\end{tikzpicture}

\caption{Growth rate vs fraction bet: concave, peaked at $f^\star=0.10$, zero near
$2f^\star$, negative beyond.}
\label{fig:kgrowth}
\end{figure}

\section{Estimation error is the bet}\label{sec:kest}

In the continuous limit $f^\star=\mu/\sigma^2$. \texttt{kelly\_est.py} perturbs the edge:
\begin{lstlisting}[numbers=none]
true: mu=0.1 sig=0.2  f*=2.50
mu off -50%:  fhat=+1.25  sizing error +0.50x
mu off -25%:  fhat=+1.88  sizing error +0.75x
mu off +0%:  fhat=+2.50  sizing error +1.00x
mu off +25%:  fhat=+3.12  sizing error +1.25x
mu off +50%:  fhat=+3.75  sizing error +1.50x
edge halves -> f* halves; bet size carries the whole estimation error
\end{lstlisting}

\begin{example}\label{ex:kest}
\textbf{Half the estimated edge, half the bet.} A 50\% error in $\hat\mu$ moves $\hat f$
by 50\%.
\begin{solution}
There is no dampening: $f^\star$ is linear in $\mu$, so every point of edge overestimation
converts directly into overbetting---into the negative side of the growth curve from
\S\ref{sec:kelly}. This is the entire case for fractional Kelly: halving the bet costs
a quarter of the growth and buys survival across the whole estimation-error range.
\end{solution}
\end{example}

\begin{figure}[htbp]\centering
\begin{tikzpicture}[font=\scriptsize]
  \draw[->,slate] (0,0)--(8.4,0) node[below left=-1pt]{$\hat\mu$ error};
  \foreach \x/\e/\f in {0.7/-50\%/1.25, 2.3/-25\%/1.88, 3.9/0/2.50, 5.5/+25\%/3.12, 7.1/+50\%/3.75}{
    \fill[accent!40] (\x,0) rectangle (\x+0.9,\f*0.7);
    \node[below] at (\x+0.45,0) {\e};
    \node[above] at (\x+0.45,\f*0.7) {\f};}
  \node[above] at (4.2,3.2) {$\hat f = \hat\mu/\sigma^2$: the whole estimation error lands in the bet};
\end{tikzpicture}

\caption{$\hat f=\hat\mu/\sigma^2$: estimation error lands whole in the position size.}
\label{fig:kest}
\end{figure}

\section{Fractional Kelly and ruin}\label{sec:frac}

\texttt{ruin\_sim.py} runs gambler's ruin on unit bets at $p{=}0.55$:
\begin{lstlisting}[numbers=none]
bankroll 5 units: ruin 33.6% of paths (barrier 3x)
bankroll 10 units: ruin 14.0% of paths (barrier 3x)
bankroll 20 units: ruin 1.5% of paths (barrier 3x)
bankroll 40 units: ruin 0.0% of paths (barrier 3x)
\end{lstlisting}

\begin{example}\label{ex:ruin}
\textbf{Bankroll depth is geometric safety.} Ruin probability falls from $33.6\%$ at
5 units to under $2\%$ at 20.
\begin{solution}
The same edge, four different outcomes---the only variable is how many bets the bankroll
can absorb. This is the quantitative face of fractional Kelly: betting a fraction of the
full stake is equivalent to running a deeper effective bankroll against the same variance.
\end{solution}
\end{example}

\begin{figure}[htbp]\centering
\begin{tikzpicture}
\begin{axis}[width=0.86\linewidth,height=0.32\linewidth,
    xlabel={bankroll (units)},ylabel={ruin prob},
    xmin=0,xmax=45,ymin=0,ymax=0.4,
    xtick={5,10,20,40},ytick={0,0.1,0.2,0.3},
    tick style={draw=none},axis line style={draw=slate},font=\scriptsize]
  \addplot[accent,very thick,mark=*] coordinates {(5,0.336)(10,0.140)(20,0.015)(40,0.0)};
  \node[font=\scriptsize,text=slate,align=center] at (axis cs:28,0.25)
    {unit bets, $p{=}0.55$\\ deeper bankroll $\Rightarrow$\\ geometric safety};
\end{axis}
\end{tikzpicture}

\caption{Ruin probability vs bankroll depth: each unit of depth buys geometric, not
linear, safety.}
\label{fig:ruin}
\end{figure}

\begin{figure}[htbp]\centering
\begin{tikzpicture}[font=\scriptsize]
  \node[draw=slate,fill=white,rounded corners=1pt,align=center,minimum width=2.5cm] at (0.5,1.0) {$f^*$\\optimal growth,\\wild path};
  \node[draw=slate,fill=white,rounded corners=1pt,align=center,minimum width=2.5cm] at (3.4,1.0) {$f^*/2$\\3/4 of the growth,\\half the drawdown};
  \node[draw=slate,fill=white,rounded corners=1pt,align=center,minimum width=2.5cm] at (6.3,1.0) {$2f^*$\\zero growth,\\certain drawdown};
  \node[text=slate] at (3.4,-0.2) {the fraction is a dial between growth and survival};
\end{tikzpicture}

\caption{The dial: half-Kelly keeps three quarters of the growth and halves the pain.}
\label{fig:halves}
\end{figure}

\section{Kelly as allocation}\label{sec:kalloc}

For multiple bets the scalar formula becomes vector Kelly, and $f^\star$ becomes the
tangency point on the return-volatility frontier. The practical constraint isn't the
math---it's that the inputs (means, covariances) are estimates with the error-carrying
property of \S\ref{sec:kest}, so real sizing layers caps on top: fraction-of-Kelly,
max-position, max-drawdown triggers.

\begin{figure}[htbp]\centering
\begin{tikzpicture}
\begin{axis}[width=0.86\linewidth,height=0.34\linewidth,
    xlabel={vol},ylabel={expected return},
    xmin=0,xmax=30,ymin=0,ymax=14,
    tick style={draw=none},axis line style={draw=slate},font=\scriptsize]
  \addplot[slate,thick,domain=0:30,samples=40] {4+0.3*x};
  \addplot[accent,thick,domain=0:30,samples=40] {x*0.45};
  \addplot[only marks,mark=*,mark size=2.5pt,red] coordinates {(20,9)};
  \node[font=\scriptsize,text=red,align=center] at (axis cs:23,8.2)
    {Kelly point\\ $\mu/\sigma^2$};
  \node[font=\scriptsize,text=slate] at (axis cs:8,10) {opportunity set};
\end{axis}
\end{tikzpicture}

\caption{Kelly picks the tangency portfolio---the point the opportunity set touches the
steepest ray.}
\label{fig:frontier}
\end{figure}

\begin{example}\label{ex:kalloc}
\textbf{Why pros bet a fraction.} Full Kelly on a true edge still delivers $\sim$98\%
median drawdowns and is optimal only for the \emph{estimated} edge, not the true one.
\begin{solution}
The practice is a stack: estimate $\mu,\sigma$, halve the bet, cap by drawdown budget,
re-check when the estimate moves. The formula supplies the scale; the layers supply
the survival.
\end{solution}
\end{example}

\begin{figure}[htbp]\centering
\begin{tikzpicture}[font=\scriptsize,every node/.style={align=center}]
  \node[draw=ink,fill=goodbg,rounded corners=1pt] (e) at (0,2.4) {edge estimate\\ $\hat\mu,\hat\sigma$};
  \node[draw=slate,fill=white,rounded corners=1pt] (h) at (-2.6,0.9) {halve it:\\ $\hat\mu/2$};
  \node[draw=slate,fill=white,rounded corners=1pt] (f) at (0.4,0.9) {$f=\hat\mu/2\hat\sigma^2$\\ (half-Kelly)};
  \node[draw=slate,fill=white,rounded corners=1pt] (c) at (3.4,0.9) {cap by\\ drawdown budget};
  \draw[->,ink] (e.south) -- (h.north);
  \draw[->,ink] (e.south) -- (f.north);
  \draw[->,ink] (e.south) -- (c.north);
  \node[text=slate] at (0.4,-0.4) {wrong on edge $\Rightarrow$ the cap, not the formula, saves you};
\end{tikzpicture}

\caption{The sizing stack: edge estimate $\to$ fractional Kelly $\to$ drawdown cap.}
\label{fig:decide}
\end{figure}

\begin{checkpoint}[Sizing audit]
\begin{itemize}
\item What does full Kelly maximize, and at what cost? (Expected log-wealth; $\sim$98\%
median drawdowns.)
\item What is $f^\star$ linear in, and why is that dangerous? ($\mu$---edge error becomes
bet-size error whole.)
\item Why is half-Kelly the common default? (3/4 of the growth, half the drawdown,
insurance against misestimation.)
\end{itemize}
\end{checkpoint}

\section{Recap}\label{sec:krecap}

Kelly is the answer to ``how much,'' and like every answer in this book it is only as
good as its inputs. The formula is exact; the edge is not. Fractional Kelly and drawdown
caps exist because the distance between $\hat\mu$ and $\mu$ is where ruin lives.

\begin{table}[htbp]\centering\small
\caption{Chapter toolbox.}
\begin{tabular}{@{}lll@{}}
\toprule
Concept & Rule & Gotcha\\
\midrule
$f^\star$ & $(bp-(1-p))/b$ & $\mu/\sigma^2$ continuous\\
Overbet & growth dies before $2f^\star$ & asymmetric penalty\\
Fractional & $f^\star/2$ default & costs 1/4 growth, buys survival\\
Ruin & geometric in bankroll depth & path risk is real\\
Caps & the formula is the scale & the cap is the survival\\
\bottomrule
\end{tabular}
\end{table}

\begin{example}[The fraction ladder]\label{ex:fracladder}
\textbf{Strategy.} Fix the 55/45 coin and sweep the betting
fraction from quarter-Kelly to $5\times$Kelly. At each step, show
the growth rate, the median final wealth over 2000 tosses, and the
median path drawdown --- the full cost curve, not just the peak.

\textbf{Solution.} \texttt{frac\_ladder.py} runs 200 seeds:
\begin{verbatim}
f=0.025: G +0.0022/toss  med final       71.9  med dd  23%  ruin<1: 0%
f=0.050: G +0.0038/toss  med final     1487.8  med dd  41%  ruin<1: 0%
f=0.100: G +0.0050/toss  med final    18325.8  med dd  70%  ruin<1: 2%
f=0.200: G -0.0001/toss  med final        0.6  med dd 100%  ruin<1: 52%
f=0.300: G -0.0162/toss  med final        0.0  med dd 100%  ruin<1: 99%
f=0.500: G -0.0889/toss  med final        0.0  med dd 100%  ruin<1: 100%
\end{verbatim}
Full Kelly ($f=0.10$) pays $18{,}326\times$ at the cost of a
median $70\%$ drawdown; double Kelly ($0.20$) already has
\emph{negative} growth and a 100\%-dd median path.

\textbf{Significance.} The growth curve is asymmetric around
$f^*$: left of it you trade return for sleep, right of it you
trade \emph{everything} for nothing. The cliff at $2f^*$ ---
where the median path visits ruin --- is why desks quote
``half-Kelly'' as the standard rather than the optimum.
\end{example}


\section{Estimation error is the actual position size}\label{sec:esterr}
Kelly's formula takes $\mu$ and $\sigma$ as \emph{given}; the
practice gives you estimates. \texttt{kelly\_est.py} perturbs
the mean on a $\mu{=}0.1,\sigma{=}0.2$ setup whose true
$f^*=2.50$:
\begin{verbatim}
true: mu=0.1 sig=0.2  f*=2.50
mu off -50%:  fhat=+1.25  sizing error +0.50x
mu off -25%:  fhat=+1.88  sizing error +0.75x
mu off +25%:  fhat=+3.12  sizing error +1.25x
mu off +50%:  fhat=+3.75  sizing error +1.50x
edge halves -> f* halves; bet size carries the whole estimation error
\end{verbatim}
\begin{example}\label{ex:esterr}
\textbf{$\hat f^*$ inherits $\hat\mu$'s error one-for-one.}
A $+50\%$ optimistic read on the edge doubles down into
$f^*=3.75$ --- the formula amplifies the mistake into the bet.
The table is linear because $f^*=\mu/\sigma^2$ is linear in
$\mu$: there is no damping.
\begin{solution}
This is the deeper reason half-Kelly persists in practice: not
superstition, but shrinkage. If $\hat\mu$ has relative error
$\epsilon$, $\hat f^*$ has the same relative error --- a
fractional haircut on $f$ is the cheapest insurance against the
overbet that ruin math punishes hardest.
\end{solution}
\end{example}
\section{Common mistakes}
\begin{itemize}
\item \textbf{Quoting Kelly without its CI.} The edge estimate carries the whole bet.
\item \textbf{Full Kelly on a real book.} $\sim$98\% drawdowns are the expected path,
  not a tail event.
\item \textbf{Treating $f^\star$ as a recommendation.} It is a scale; the fraction is
  a separate decision.
\item \textbf{Ignoring correlated bets.} Vector Kelly without the covariance is a
  different (larger) bet.
\item \textbf{``Growth-optimal'' framing for clients.} It optimizes log wealth, not
  any human's utility.
\item \textbf{Recomputing $f^\star$ per trade on stale $\hat\mu$.} The input estimate is
  the strategy.
\item \textbf{Optimising the peak, ignoring the path.}
Full Kelly pays best \emph{median} growth but a 70\% median
drawdown (Example~\ref{ex:fracladder}) --- an investor who cannot
hold the drawdown never collects the median.
\item \textbf{``A bit more can't hurt.''} Growth goes negative at
$2f^*$ --- the cost of overbetting is not a flatter curve but a
cliff.\end{itemize}

\section{Exercises}
\begin{exercise}[easy]
Write the Kelly fraction for a $p{=}0.6$, $b{=}1$ coin.
\end{exercise}
\begin{solution}
$f^\star=p-(1-p)=0.2$.
\end{solution}

\begin{exercise}[easy]
What does $2f^\star$ do to growth on the even-money coin?
\end{exercise}
\begin{solution}
Drives it to approximately zero---the overbet penalty is asymmetric.
\end{solution}

\begin{exercise}[easy]
What is the continuous-limit $f^\star$?
\end{exercise}
\begin{solution}
$f^\star=\mu/\sigma^2$.
\end{solution}

\begin{exercise}[easy]
Ruin probability at 20-unit bankroll, $p{=}0.55$ unit bets?
\end{exercise}
\begin{solution}
About $1.5\%$---vs $33.6\%$ at 5 units.
\end{solution}

\begin{exercise}[easy]
What does half-Kelly cost in growth?
\end{exercise}
\begin{solution}
Roughly a quarter of the full-Kelly growth rate---the standard survival insurance.
\end{solution}

\begin{exercise}
Derive $f^\star$ by maximizing $g(f)$ on the even-money coin.
\end{exercise}
\begin{solution}
$g'(f)=pb/(1+fb)-(1-p)/(1-f)=0 \Rightarrow f=(pb-(1-p))/(b+p b-f b)$... solving gives
$f^\star=(bp+p-1)/(b)$... for $b{=}1$: $f^\star=2p-1$.
\end{solution}

\begin{exercise}
A 20\%-edge strategy has $\sigma{=}25\%$. Full and half Kelly?
\end{exercise}
\begin{solution}
$f^\star=0.20/0.0625=3.2$; half-Kelly $1.6$---numbers that require leverage, which is
why the cap layer exists.
\end{solution}

\begin{exercise}
Show that a 25\% overestimate of $\mu$ costs more than growth---where does the error land?
\end{exercise}
\begin{solution}
$\hat f$ is $1.25\times$ true $f^\star$: on the growth curve you are right of the peak,
where the derivative is negative.
\end{solution}

\begin{exercise}
Why is ruin probability geometric in bankroll, not linear?
\end{exercise}
\begin{solution}
Each unit of depth requires an independent losing streak; probabilities multiply,
so risk falls as $(q/p)^N$.
\end{solution}

\begin{exercise}
What does vector Kelly require that scalar Kelly does not?
\end{exercise}
\begin{solution}
The covariance matrix---correlated bets share variance, so the tangency solution is
different from independent scalars.
\end{solution}

\begin{exercise}
Why does half-Kelly dominate full Kelly under estimation error?
\end{exercise}
\begin{solution}
Overbetting is penalized asymmetrically; halving keeps you on the left (growth-positive)
side of the peak under any plausible $\hat\mu$ error.
\end{solution}

\begin{exercise}
A drawdown budget of 20\%---how does it translate to a fraction cap?
\end{exercise}
\begin{solution}
Find $f$ with $P(\text{dd}>20\%)$ below tolerance via sim or the ruin formula; the cap
is whichever of $\{f^\star, f_{\text{cap}}\}$ is smaller.
\end{solution}

\begin{exercise}[hard]
Derive the growth cost of half-Kelly exactly for the 55/45 coin, and verify the
``quarter of growth'' heuristic.
\end{exercise}
\begin{solution}
$g(f^\star)=0.0050$, $g(f^\star/2)\approx0.0038$---ratio $\approx0.76$: three quarters
of growth for half the drawdown, confirmed.
\end{solution}

\begin{exercise}[hard]
Explain why ``Kelly is optimal'' and ``Kelly ruins you'' are both true, naming the two
senses of optimality.
\end{exercise}
\begin{solution}
Growth-optimal (max asymptotic log wealth) vs path-wise survival---the maximizing
trajectory passes through arbitrary drawdowns; optimality is a statement about $t\to\infty$,
not about staying in the game.
\end{solution}

\begin{exercise}[hard]
Construct a sizing stack for a strategy with $\hat\mu/\sigma^2=2$ and an estimated
$N{=}300$ trial count---what does each layer do?
\end{exercise}
\begin{solution}
Estimate $\hat\mu$; apply $N$-uncertainty (CI on $\mu$); take fractional $f=1$;
apply drawdown cap $\to$ final fraction. Each layer converts a different error source
into a smaller bet.
\end{solution}

\begin{exercise}[easy]
From the ladder: growth at $f=0.05$ is $+0.0038$ and at $f=0.20$
is $-0.0001$. Which is $f^*$ and which is $2f^*$?
\end{exercise}
\begin{solution}
$f^*=0.10$ sits between them: half-Kelly still grows
($+0.0038$), double-Kelly has already crossed zero growth ---
the signature of the cliff.
\end{solution}
\begin{exercise}[medium]
Your investor will quit at a 50\% drawdown. Which fraction from
the ladder respects the constraint, and what growth does it
cost?
\end{exercise}
\begin{solution}
$f\approx0.05$ (median dd $41\%<50\%$): you give up
$\sim0.0012$/toss of growth vs $f^*$ to stay inside the mandate.
\end{solution}
\begin{exercise}[hard]
Show why the ruin column jumps to $\sim$50\% exactly at $2f^*$:
prove $G(2f^*)\le0$ for the even-payoff coin, i.e.\
$p\ln(1+2f^*)+(1-p)\ln(1-2f^*)\le0$.
\end{exercise}
\begin{solution}
$f^*=2p-1$: at $2f^*$, $G=(1-p)\ln(1-2(2p-1))+p\ln(1+2(2p-1))$.
For $p=0.55$: $0.45\ln0.8+0.55\ln1.2 = -0.1004+0.1003\approx0$
--- the cliff edge is algebraic, not a simulation artifact.
\end{solution}

\section*{Notes}
Kelly is the only sizing rule with a theorem behind it, and the theorem assumes the edge
is known exactly---which it never is. The practitioner's answer is always a fraction,
always inside caps, and always subordinate to the estimation error from ch12: the bet
you cannot afford is the one whose edge you measured wrong.

\section*{Further reading}
\begin{itemize}
\item J.~L.~Kelly, ``A New Interpretation of Information Rate'' --- the 1956 paper:

\url{https://archive.org/details/bstj35-4-917}
\item E.~O.~Thorp, ``The Kelly Criterion in Blackjack, Sports Betting, and the Stock Market'':

\url{https://wayback.archive-it.org/all/20160815092123/http://www.edwardothorp.com/sitebuildercontent/sitebuilderfiles/Kelly_Criterion_in_Blackjack_Sports_Betting_and_the_Stock_Market.pdf}
\item MacLean, Thorp \& Ziemba, ``The Kelly Capital Growth Investment Criterion'':

\url{https://www.worldscientific.com/worldscibooks/10.1142/7598}
\item L.~de Prado lectures --- position sizing:

\url{https://www.quantresearch.org/Lectures.htm}
\end{itemize}

\begin{rulebox}[Watch instead]
\begin{itemize}
\item ``12. Risk Analysis'' --- MIT 18.S096 (80:32). The risk side of sizing; re-verified
2026-10-11.
\url{https://www.youtube.com/watch?v=qlytPllimpQ}
\item ``1. Introduction and Financial Terms'' --- MIT 18.S096 (60:30). Utility framing;
re-verified 2026-10-11.
\url{https://www.youtube.com/watch?v=wvXDB9dMdEo}
\end{itemize}
\end{rulebox}
